% Lecture 11: four open questions; optional later review, no answer gate.
\documentclass[a4paper,12pt]{article}
\usepackage[margin=19mm]{geometry}
\usepackage{amsmath,amssymb,graphicx,xcolor}
\usepackage{fontspec}
\usepackage{xepersian}
\settextfont{XB Roya}
\setlatintextfont{Latin Modern Roman}
\definecolor{ink}{HTML}{173D75}
\graphicspath{{images/}}
\setlength{\parindent}{0pt}
\setlength{\parskip}{5pt}
\linespread{1.12}
\newcommand{\qtitle}[2]{{\Large\bfseries\color{ink} پرسش \lr{#1}: #2}\par}
\newcommand{\answer}{\par\medskip\textbf{پاسخ یا صورت مورد تأیید و تاریخ:}\par\vspace{18mm}\hrule\medskip}
\newcommand{\sourcepage}[1]{%
  \textbf{نشانی منبع:} نسخهٔ \lr{primary}، صفحهٔ #1.\par
  \begin{LTR}\footnotesize
  \texttt{sources/lectures/11 Notherian and Artinian modules (2).pdf}
  \end{LTR}}
\newcommand{\pending}{\textbf{وضعیت:} باز، برای بررسی اختیاریِ بعدی؛ تا تاریخ
  \lr{2026-10-25} پاسخی ثبت نشده است.\par}
\newcommand{\evidence}[2][65mm]{\begin{center}
  \includegraphics[width=\linewidth,height=#1,keepaspectratio]{#2}
  \end{center}}
\begin{document}
{\LARGE\bfseries پرسش‌های ابهام جلسهٔ ۱۱}\par
{\large جزوهٔ جبر پیشرفتهٔ دکتر خسروی}\par
چهار پرسش زیر به اثبات‌ها و فرض‌های منبع مربوط‌اند. شمارهٔ صفحه‌ها، شمارهٔ
یک‌مبنای صفحه در فایل پی‌دی‌اف است. صورت‌های موقت و موارد موکول‌شده در اسلایدها
مشخص‌اند. این پرونده برای بررسی بعدی است و پاسخ، شرط پایان این دور رونویسی نیست.
\medskip\hrule\medskip

\qtitle{L11-Q01}{انتخاب اشتراک‌های متناهی در خانوادهٔ دلخواه}
\pending
\sourcepage{۱، نیمهٔ پایین، اثبات آرتینی بودن به هم‌تولید متناهی خارج‌قسمت‌ها}
\textbf{عبارت منبع:} برای خانوادهٔ $\{M_i\}_{i\in I}$ با اشتراک $N$، زنجیرهٔ
$M_1\supsetneq M_1\cap M_2\supsetneq\cdots$ نوشته و سپس نتیجه گرفته شده است:
\[
  \exists k:\quad \bigcap_{i=1}^{k}M_i=\bigcap_{i\in I}M_i=N.
\]
\evidence[52mm]{L11-Q01-intersections.png}
\textbf{خوانش‌های ممکن:} $I$ قابل شمارش فرض شده است؛ یا منظور انتخاب پی‌درپی
اعضایی از خانواده است که تا پیش از رسیدن به $N$، اشتراک را اکیداً کوچک کنند.

\textbf{دلیل پرسش:} توقف اشتراک‌های یک دنبالهٔ دلخواه به‌تنهایی ثابت نمی‌کند
اشتراک آن برابر اشتراک کل خانوادهٔ دلخواه است؛ نحوهٔ انتخاب در جزوه تصریح نشده است.

\textbf{وضعیت در اسلاید:} آغاز اثبات حفظ شده و گام انتخاب اشتراک‌ها صریحاً
به بررسی بعدی موکول شده است؛ تکمیل این گام به‌عنوان متن جزوه عرضه نشده است.

\textbf{پرسش:} آیا منظور همان انتخاب پی‌درپیِ اشتراک‌های اکیداً کوچک‌تر است
و اجازه می‌دهید این مرحله در اثبات اسلاید صریح شود؟ یا فرض دیگری برای $I$ منظور است؟
\answer
\newpage

\qtitle{L11-Q02}{شمول سره در زنجیره‌های القایی}
\pending
\sourcepage{۴، نیمهٔ پایین، و ۵، سطر اول؛ اثبات جهت عکس معیار آرتینی}
\textbf{عبارت منبع:} از $M_1\supsetneq M_2\supsetneq\cdots$، زنجیره‌های
$M_i+N$، $(M_i+N)/N$ و $M_i\cap N$ هم با شمول سره نوشته شده‌اند.
\evidence[80mm]{L11-Q02-images-chain.png}
\textbf{خوانش‌های ممکن:} حفظ شمول سره به‌صورت لفظی؛ یا اصلاح آن به شمول
غیرسره در سه زنجیرهٔ القایی.

\textbf{دلیل پرسش:} جمع با $N$، تصویر در خارج‌قسمت و اشتراک با $N$ لزوماً
شمول سره را حفظ نمی‌کنند. برای استفاده از شرط زنجیر نزولی، شمول غیرسره کافی است.

\textbf{انتخاب موقت در اسلاید:} شمول غیرسره به‌طور آشکار به‌عنوان اصلاح
موقت استفاده شده است؛ این انتخاب، پاسخ یا تأیید کاربر محسوب نمی‌شود.

\textbf{پرسش:} آیا در سه زنجیرهٔ القایی، همهٔ علامت‌های $\supsetneq$ را به
$\supseteq$ تغییر دهم؟ زنجیرهٔ فرض خلافِ اولیه همچنان می‌تواند سره بماند.
\answer
\newpage

\qtitle{L11-Q03}{نوتری یا آرتینی در توقف زنجیرهٔ نزولی؟}
\pending
\sourcepage{۴، پایین صفحه، و ۵، بالای صفحه، پس از زنجیرهٔ اشتراک‌ها}
\textbf{عبارت منبع:} در صفحهٔ ۴ «چون $M/N$ نوتری» و در صفحهٔ ۵ «چون $N$
نوتری» خوانده می‌شود؛ سپس توقف زنجیره‌های نزولی نتیجه گرفته شده است.
\evidence[52mm]{L11-Q03-quotient-condition.png}
\evidence[34mm]{L11-Q03-condition.png}
\textbf{خوانش‌های ممکن:} «نوتری» مطابق ظاهر دست‌نویس؛ یا «آرتینی» مطابق
فرض قضیه و نوع زنجیره.

\textbf{دلیل پرسش:} در صورت قضیه $N$ و $M/N$ آرتینی‌اند. نوتری بودن به‌تنهایی
توقف زنجیرهٔ نزولی را نتیجه نمی‌دهد.

\textbf{انتخاب موقت در اسلاید:} در هر دو محل از فرض آرتینی بودن استفاده
شده و موقت بودنِ این اصلاح در تمام ادامهٔ اثبات مشخص است.

\textbf{پرسش:} آیا هر دو عبارت را به «چون $M/N$ آرتینی است» و «چون $N$
آرتینی است» اصلاح کنم؟
\answer
\newpage

\qtitle{L11-Q04}{فرض یکانی بودن در دو قضیهٔ پایانی}
\pending
\sourcepage{۸، قضیهٔ میانی و قضیه و اثبات پایین صفحه}
\textbf{عبارت منبع:} $R$ یکدار است و «هر $R$-مدول با تولید متناهی» مطرح
می‌شود؛ در قضیهٔ آخر نیز $M\ne0$ یک $R$-مدول «دلخواه» است. قید یکانی صریح نیست.
\evidence[45mm]{L11-Q04-fg-statement.png}
\evidence[43mm]{L11-Q04-cyclic-proof.png}
\textbf{خوانش‌های ممکن:} همهٔ مدول‌ها در این دو قضیه از ابتدا یکانی فرض
شده‌اند؛ یا عبارت «دلخواه» مدول‌های غیریکانی را هم شامل می‌شود.

\textbf{دلیل پرسش:} پوشایی $R^n\to M$ در اثبات اول و ناصفر بودن $Rx$ برای
$x\ne0$ در اثبات آخر به این قرارداد وابسته‌اند. در جزوهٔ درس، یکانی بودن
مدول پیش‌تر به‌صورت فرضی جداگانه تعریف شده است.

\textbf{انتخاب موقت در اسلاید:} مدول‌های چپ یکانی فرض شده‌اند؛ این
فرض افزوده، کنار هر دو قضیه و اثبات مربوط آشکارا موقت اعلام شده است.

\textbf{پرسش:} آیا در هر دو قضیه و اثبات مربوط، «$R$-مدول چپ یکانی» را
صریح بنویسم؟ در غیر این صورت، قرارداد یا صورت دقیق مورد نظر را مشخص کنید.
\answer
\end{document}
